% Options for packages loaded elsewhere
\PassOptionsToPackage{unicode}{hyperref}
\PassOptionsToPackage{hyphens}{url}
\documentclass[
]{article}
\usepackage{xcolor}
\usepackage{amsmath,amssymb}
\setcounter{secnumdepth}{-\maxdimen} % remove section numbering
\usepackage{iftex}
\ifPDFTeX
  \usepackage[T1]{fontenc}
  \usepackage[utf8]{inputenc}
  \usepackage{textcomp} % provide euro and other symbols
\else % if luatex or xetex
  \usepackage{unicode-math} % this also loads fontspec
  \defaultfontfeatures{Scale=MatchLowercase}
  \defaultfontfeatures[\rmfamily]{Ligatures=TeX,Scale=1}
\fi
\usepackage{lmodern}
\ifPDFTeX\else
  % xetex/luatex font selection
\fi
% Use upquote if available, for straight quotes in verbatim environments
\IfFileExists{upquote.sty}{\usepackage{upquote}}{}
\IfFileExists{microtype.sty}{% use microtype if available
  \usepackage[]{microtype}
  \UseMicrotypeSet[protrusion]{basicmath} % disable protrusion for tt fonts
}{}
\makeatletter
\@ifundefined{KOMAClassName}{% if non-KOMA class
  \IfFileExists{parskip.sty}{%
    \usepackage{parskip}
  }{% else
    \setlength{\parindent}{0pt}
    \setlength{\parskip}{6pt plus 2pt minus 1pt}}
}{% if KOMA class
  \KOMAoptions{parskip=half}}
\makeatother
\setlength{\emergencystretch}{3em} % prevent overfull lines
\providecommand{\tightlist}{%
  \setlength{\itemsep}{0pt}\setlength{\parskip}{0pt}}
\usepackage{bookmark}
\IfFileExists{xurl.sty}{\usepackage{xurl}}{} % add URL line breaks if available
\urlstyle{same}
\hypersetup{
  pdftitle={02\_04\_Tag\_1\_Docker\_SQL\_Infrastruktur},
  hidelinks,
  pdfcreator={LaTeX via pandoc}}

\title{02\_04\_Tag\_1\_Docker\_SQL\_Infrastruktur}
\author{}
\date{}

\begin{document}
\maketitle

\textbf{Datum:} 09.03.2026\\
\textbf{Bearbeiter:} Samuel

\subsection{🎯 Tagesziele}\label{tagesziele}

Docker-Umgebung auf Ubuntu-Server verifizieren

BetaTrade-Kundendatenbank via Docker-Compose bereitstellen

Business Intelligence Abfragen (SQL) zur Kundenstruktur durchführen

Backup-Prozess für MySQL-Container validieren

\subsection{👥 Gruppenarbeit}\label{gruppenarbeit}

\begin{enumerate}
\tightlist
\item
  \textbf{Infrastruktur-Check:} Gemeinsame Durchsicht der
  \texttt{docker-compose.yml} (Mapping von Ports und Volumes).
\item
  \textbf{SQL-Review:} Besprechung der JOIN-Logik für die Verknüpfung
  von \texttt{customers}, \texttt{accounts} und \texttt{trades}.
\end{enumerate}

\subsection{👤 Einzelarbeit}\label{einzelarbeit}

\begin{enumerate}
\tightlist
\item
  \textbf{Container-Management:} Implementierung der Lifecycle-Befehle
  (Up/Down/Logs).
\item
  \textbf{Datenanalyse:} Erstellung und Test der SQL-Abfragen für die
  Neukunden-Analyse 2024.
\item
  \textbf{Dokumentation:} Erstellung des Technischen Guides für Docker
  \& SQL (siehe Analysen).
\end{enumerate}

\subsection{📝 Reflexion / Offene Punkte}\label{reflexion-offene-punkte}

\begin{itemize}
\tightlist
\item
  Die Container-Lösung bietet eine saubere Trennung von Anwendung und
  Datenbank.
\item
  Herausforderung: SQL-Dumps müssen automatisiert in den
  Cloud-Backup-Prozess integriert werden.
\item
  Nächster Schritt: Integration der Datenbank in das Monitoring
  (Themenfeld 4).
\end{itemize}

\end{document}
